\documentclass[aspectratio=169,11pt]{beamer}

\usepackage[utf8]{inputenc}
\usepackage{amsmath,amssymb,amsfonts}
\usepackage{physics}          % for \dv, \pdv, etc.
\usepackage{mathtools}
\usepackage{bm}
\usepackage{tikz}
\usepackage{booktabs}
\usepackage{graphicx}

\usetheme{Madrid}
\usecolortheme{default}

% Cleaner math
\setlength{\abovedisplayskip}{6pt}
\setlength{\belowdisplayskip}{6pt}
\setlength{\abovedisplayshortskip}{4pt}
\setlength{\belowdisplayshortskip}{4pt}

\title{What is Velocity?}
\subtitle{A Journey Through Physics}
\author{}
\date{}

\begin{document}

%================================================
\begin{frame}
\titlepage
\end{frame}

%================================================
\begin{frame}{Introduction}
Let's define velocity!
\end{frame}

%================================================
\begin{frame}[fragile]{Introduction (cont.)}
{\footnotesize
\begin{align*}
V &\equiv (Q_{\xi},X) = \mathfrak{X}_{Q_{\xi}}[X] \\[0.6ex]
\delta[Q_{\xi}] &= K(\delta,\mathfrak{X}_{Q_{\xi}}) = \Omega_{\mathrm{ext}}(\delta,\mathfrak{X}_{Q_{\xi}}) - G(\delta,\mathfrak{X}_{Q_{\xi}}) \\[0.8ex]
&= \Omega(\delta,\mathfrak{L}_{\xi}) + Q_{[\mathbb{X}(\delta),\mathbb{X}(\mathfrak{L}_{\xi})]}
  + \oint_{\partial\Sigma} \mathbb{X}(\mathfrak{L}_{\xi})\mathbin{\lrcorner}\vartheta(\delta) \\
&\qquad + \oint_{\partial\Sigma} \mathbb{X}(\mathfrak{L}_{\xi})\mathbin{\lrcorner}\vartheta(\delta) - G(\delta,\mathfrak{L}_{\xi}) \\[0.8ex]
&= \Omega_{\mathrm{ext}}(\delta,\mathfrak{L}_{\xi}) - G(\delta,\mathfrak{L}_{\xi})
\end{align*}
}
\end{frame}

%================================================
\begin{frame}[fragile]{Introduction (cont.)}
{\footnotesize
\begin{align*}
&= \delta\!\left[\int_{\Sigma}\!\bigl(\vartheta(\mathfrak{L}_{\xi}) - \xi\mathbin{\lrcorner}L\bigr)\right]
   - Q_{\delta[\xi]} - \oint_{\partial\Sigma}\xi\mathbin{\lrcorner}\vartheta(\delta) \\
&\qquad + Q_{[\mathbb{X}(\delta),\mathbb{X}(\mathfrak{L}_{\xi})]}
  + \oint_{\partial\Sigma}\mathbb{X}(\mathfrak{L}_{\xi})\mathbin{\lrcorner}\vartheta(\delta) \\[1ex]
&= \delta\!\left[\int_{\Sigma}\!\bigl(\vartheta(2\nabla_{(a}\xi_{b)},\mathcal{L}_{\xi}\psi^{I}) - \xi\mathbin{\lrcorner}L\bigr)\right]
   - Q_{\delta[\xi]} \\
&\qquad + Q_{[\mathbb{X}(\delta),\xi]} + \oint_{\partial\Sigma}\xi\mathbin{\lrcorner}\vartheta(\delta)
\end{align*}
}
\end{frame}

%================================================
\begin{frame}[fragile]{Introduction (cont.)}
{\footnotesize
\begin{align*}
\dot{F} &= (Q_{\xi},F) \\[0.6ex]
(Q_{\xi},g_{ab}) &= \mathcal{L}_{\xi}g_{ab} = 2\nabla_{(a}\xi_{b)} \\
(Q_{\xi},\psi^{I}) &= \mathcal{L}_{\xi}\psi^{I} \\
(Q_{\xi},x^{\mu}) &= \xi^{\mu} \\[0.8ex]
\Gamma^{c}{}_{ab} &= \tfrac12 g^{cd}(\partial_{a}g_{db}+\partial_{b}g_{da}-\partial_{d}g_{ab}) \\
g^{ac}g_{cb}&=\delta^{a}{}_{b},\qquad\nabla_{c}g_{ab}=0 \\[0.8ex]
(Q_{\xi},Q_{\xi}) &= -G(\mathfrak{L}_{\xi},\mathfrak{L}_{\xi})
  = \oint_{\partial\Sigma}\xi\mathbin{\lrcorner}\vartheta(\mathfrak{L}_{\xi})
\end{align*}
}

\vspace{0.8em}
\centering\large \ldots uhm, well\ldots maybe I am wrong about that.
\end{frame}

%================================================
\begin{frame}{Basic Definition}


Velocity can generally be defined as the change in distance traveled divided by the time it takes to travel:
\[
v = \frac{\Delta x}{\Delta t}
\]

Taking the infinitesimal limit:
\[
v = \dv{x}{t}
\]

Velocity also contains direction. The easiest way to represent it is as a vector:
\[
\mathbf{\dot{r}} = \dv{\mathbf{r}}{t}
= \Bigl\langle \dv{x}{t},\;\dv{y}{t},\;\dv{z}{t}\Bigr\rangle
\]
\end{frame}

%================================================
\begin{frame}{Newtonian Mechanics}
Newton's second law:
\[
\mathbf{F} = m\,\dv[2]{\mathbf{r}}{t}
\]

Solving for velocity:
\[
\dot{r} = \int\frac{\mathbf{F}}{m}\,dt
\]

Keeping track of the reference frame is simple:
\[
\dot{r}' = \int\frac{\mathbf{F}}{m}\,dt - \dot{r}_{\text{reference frame}}
\]
\end{frame}

%================================================
\begin{frame}{Newtonian Mechanics (cont.)}
Relying on frames of reference creates problems (e.g.\ fictitious centrifugal forces).

One way to make Newtonian mechanics generally covariant is the \textbf{Newton--Cartan} approach:
\begin{itemize}
  \item Space-time structure similar to general relativity
  \item But time remains absolute
\end{itemize}

Velocity is then defined as
\[
V^{\mu} = \dv{x^{\mu}}{t} = \bigl(1,\;\dv{x^{i}}{t}\bigr)
\]

While this looks simple, it is an extreme simplification.
\end{frame}

%================================================
\begin{frame}{Special Relativity}
In special relativity, time becomes another relative dimension.

The velocity vector becomes four-dimensional:
\[
U^{\mu} = \dv{x^{\mu}}{\tau}
\]
where $\tau$ is \textbf{proper time} (the time experienced by the traveler).

Galilean relative velocity
\[
v_{\mathrm{rel}} = v_{1} - v_{0}
\]
gains relativistic corrections:
\[
v_{\mathrm{rel}} = \frac{v_{1}-v_{0}}{1-v_{1}v_{0}/c^{2}}
\]
\end{frame}

%================================================
\begin{frame}{General Relativity}
Absolute acceleration creates paradoxes and preferential frames.

Combining the covariant Newtonian framework with Einsteinian proper time yields the same-looking equation
\[
U^{\mu} = \dv{x^{\mu}}{\tau}
\]
but now with the normalization
\[
g_{\mu\nu}U^{\mu}U^{\nu} = -c^{2}.
\]

The geometry of spacetime is curved, further separating observers. The rest proceeds as in special relativity, but with additional corrections expressible in the language of Riemannian geometry and geodesics.
\end{frame}

%================================================
\begin{frame}{Warp Drives}
Warp drives \emph{move} faster-than-light motion.

\begin{itemize}
  \item Most solutions have the matter itself not moving at all.
  \item Spacetime contracts in front and expands behind the ship.
  \item Proper velocity of the ship is essentially $0$, yet an external observer sees it traverse great distances in short coordinate time.
\end{itemize}

This makes the definition of ``velocity'' rather challenging to define.
\end{frame}

%================================================
\begin{frame}{Dynamics --- Phase Space}

In phase space the fundamental coordinates are generalized positions $q^{i}$ and momenta $p_{i}$.

If the flow is generated by a Hamiltonian $H$, then
\[
v = \dot{q}^{i} = \pdv{H}{p_{i}}.
\]

Velocity is simply one component of the Hamiltonian vector field on phase space.
\end{frame}

%================================================
\begin{frame}{Poisson Structure}
We introduce the Poisson bracket
\[
\dot{f} = \{f,H\}
= \pdv{f}{q^{i}}\pdv{H}{p_{i}} - \pdv{f}{p_{i}}\pdv{H}{q^{i}}.
\]

The entire vector field can then be written
\[
X_{H} = \{\,\cdot\,,H\}.
\]

The underlying geometry is symplectic (we will not go into the details here).
\end{frame}

%================================================
\begin{frame}{Other Approaches (Speed Run)}
\begin{itemize}
  \item \textbf{Fluid dynamics}: a vector field that assigns a Newtonian velocity to every point in space.
  \item \textbf{Quantum mechanics}: velocity becomes a probability current
  \[
  \mathbf{v} = \frac{\mathbf{j}}{|\psi|^{2}}.
  \]
  \item \textbf{Kinetic theory}: combines the probabilistic nature of QM with the continuum of fluids.
  \item \textbf{Geometric mechanics}: velocity is a ``natural element of a tangent bundle.''
\end{itemize}
\end{frame}

%================================================

%================================================
\begin{frame}
\begin{center}
\Huge Thank you
\end{center}
\end{frame}

\end{document}
